\section{The semigroup structure of the greatest ambit}\label{sec:semigroup}

The universal property of the greatest ambit turns it into a compact
semigroup. Minimal flows inside it are minimal ideals, and comparing it with
the enveloping semigroup of the universal minimal flow leads to Ellis's
problem and to a conjecture of Pestov. Some sources, such as
\cite{BartosovaZucker2019}, work with right actions; their semigroups are then
left topological rather than right topological.

\begin{definition}[Right topological semigroups]\label{def:rt-semigroup}
A \emph{compact right topological semigroup} is a semigroup $S$ with a compact
topology in which, for every $q\in S$, the map $p\mapsto pq$ is continuous. An
element $u\in S$ is an \emph{idempotent} if $uu=u$. A \emph{left ideal} is a
nonempty set $I\subseteq S$ with $SI\subseteq I$. Terminology varies: such
semigroups are called left topological in \cite{Uspenskij2000}.
\end{definition}

\begin{proposition}\label{prop:ambit-semigroup}
\uses{def:rt-semigroup, prop:greatest-ambit}
For $x\in S(G)$, let $\rho_x\colon(S(G),1)\to(\overline{Gx},x)$ be the morphism
of ambits given by Proposition~\ref{prop:greatest-ambit}, and put
$yx=\rho_x(y)$. This makes $S(G)$ a compact right topological semigroup with
identity $1$, in which left multiplication by $g\in G$ is the action of $g$
\cite[Theorem~3.1]{Uspenskij2000}.
\end{proposition}

\begin{proof}
Right multiplication by $x$ is $\rho_x$, which is continuous. Since $\rho_x$ is
a $G$-map with $\rho_x(1)=x$, we get $gx=\rho_x(g)$ for $g\in G$; also
$\rho_1$ is the identity, so $y1=y$. For associativity, $\rho_z\circ\rho_x$ is
a $G$-map sending $1$ to $xz$, hence a morphism of ambits
$(S(G),1)\to(\overline{Gxz},xz)$, so it equals $\rho_{xz}$ by uniqueness;
that is, $(yx)z=y(xz)$.
\end{proof}

\begin{proposition}\label{prop:minimal-ideals}
\uses{prop:ambit-semigroup, cor:umf-in-ambit}
\begin{enumerate}
\item The subflows of $S(G)$ are exactly its closed left ideals, and its
  minimal subflows are exactly its minimal left ideals. Each of them is
  isomorphic to $M(G)$.
\item Every minimal left ideal $M$ contains an idempotent $u$, and
  $p\mapsto pu$ is a $G$-map $S(G)\to M$ that restricts to the identity on
  $M$. In particular $M(G)$ is a retract of $S(G)$.
\end{enumerate}
\end{proposition}

\begin{proof}
(1) A left ideal is invariant under left multiplication by $G$. Conversely, if
$Y$ is a subflow, $x\in Y$ and $y=\lim g_i$ with $g_i\in G$, then
$yx=\lim g_ix\in Y$, by continuity of $\rho_x$. Every left ideal contains a
left ideal $S(G)x=\rho_x(S(G))$, which is closed, so minimal left ideals are
closed and are the minimal subflows. They are isomorphic to $M(G)$ by
Corollary~\ref{cor:umf-in-ambit}.

(2) A minimal left ideal $M$ is a closed subsemigroup, and every compact right
topological semigroup contains an idempotent, by a theorem of Ellis; see
\cite[Theorem~3.2]{Uspenskij2000} for a short proof. Let $u\in M$ be an idempotent. Then
$S(G)u$ is a left ideal contained in $M$, so $M=S(G)u$, and $pu=p$ for every
$p=qu\in M$. The map $p\mapsto pu$ is continuous, and it is a $G$-map by
associativity.
\end{proof}

\begin{definition}[Enveloping semigroup]\label{def:enveloping}
\uses{def:flow, def:ambit}
The \emph{enveloping semigroup} $E(X)$ of a $G$-flow $X$ is the closure, in
the product space $X^X$, of the set of maps $x\mapsto gx$ for $g\in G$. It is a
compact right topological semigroup under composition, and $G$ acts on it by
$(gp)(x)=g\,p(x)$; with the identity map as base point, it is a $G$-ambit. By
Proposition~\ref{prop:greatest-ambit} there is a canonical morphism of ambits
$S(G)\to E(M(G))$.
\end{definition}

\begin{remark}\label{rem:ellis-question-semigroup}
\uses{def:enveloping, rem:ellis-problem, prop:ea-trivial-umf}
Ellis's problem (Remark~\ref{rem:ellis-problem}) asks whether the canonical
map $S(G)\to E(M(G))$ is always an isomorphism. For a nontrivial extremely
amenable group it is not, since $E(M(G))$ is a point while $S(G)$ is not
\cite{Pestov1998}.
\end{remark}

\begin{theorem}[Glasner \cite{Glasner1998}]\label{thm:glasner-monothetic}
\uses{rem:ellis-question-semigroup, def:ea}
\begin{enumerate}
\item There is an infinite monothetic Polish group (one with a dense cyclic
  subgroup) that is extremely amenable. This answers a question of Mitchell.
\item For $\mathbb{Z}$ with the discrete topology, the greatest ambit
  $\beta\mathbb{Z}$ is not isomorphic to $E(M(\mathbb{Z}))$.
\end{enumerate}
\end{theorem}

\begin{conjecture}[Pestov]\label{conj:pestov}
\uses{rem:ellis-question-semigroup}
For a topological group $G$, the canonical map $S(G)\to E(M(G))$ is an
isomorphism if and only if $G$ is precompact, that is, isomorphic to a
subgroup of a compact group \cite[page~4163]{Pestov1998}.
\end{conjecture}

\begin{remark}\label{rem:pestov-conjecture}
\uses{conj:pestov, ex:umf-compact}
If $G$ is precompact, $S(G)$, $M(G)$ and $E(M(G))$ are all isomorphic to the
completion of $G$, which is a compact group \cite{BartosovaZucker2019}; compare
Example~\ref{ex:umf-compact}. So the conjecture concerns non-precompact
groups. Pestov showed that the canonical map is an isomorphism if and only if
the retractions of $S(G)$ onto a fixed minimal subflow, as in
Proposition~\ref{prop:minimal-ideals}, separate the points of $S(G)$
\cite{BartosovaZucker2019}.
\end{remark}

\begin{theorem}[Glasner--Tsankov--Weiss--Zucker \cite{GTWZ2021}]\label{thm:gtwz}
\uses{conj:pestov}
Pestov's conjecture holds for discrete groups: if $G$ is an infinite discrete
group, then $\beta G$ is not isomorphic to $E(M(G))$.
\end{theorem}

\begin{theorem}[Bartošová--Zucker \cite{BartosovaZucker2019}]\label{thm:bartosova-zucker}
\uses{conj:pestov, ex:kpt-computations}
Let $G$ be $S_\infty$, the automorphism group of a random $K_n$-free graph, or
the automorphism group of a random $r$-uniform hypergraph. Then $S(G)$ is not
isomorphic to $E(M(G))$.
\end{theorem}

\begin{remark}\label{rem:bartosova-zucker}
\uses{thm:bartosova-zucker, thm:gtwz}
The proof uses the explicit universal minimal flows of
Section~\ref{sec:kpt-umf} and only works when $M(G)$ is metrizable. The
methods of \cite{GTWZ2021} are of a different kind: they rely on a
disjointness property that groups with metrizable universal minimal flow
never have \cite{BartosovaZucker2019}.
\end{remark}
